\Problem{Números d-espaiats}


\Statement

Donada una base $b \geq 2$ i una xifra $0 \leq d \leq b$, diem que un enter no negatiu $n$ és \emph{d-espaiat} si $n$ expressat en base $b$ no té cap parell de xifres consecutives tals que totes dues són diferents de $d$. Per exemple:

\begin{itemize}
\item En base $10$, l'enter $145447$ és $4$-espaiat, mentre que el $145$ és $4$-espaiat però no és $1$-espaiat (el parell de dígits $45$ no compleix la condició).

\item En base $2$, l'enter $9$ ($1001$ en binari) és $0$-espaiat però no $1$-espaiat (el parell de xifres $00$ no compleix la condició).

\item En base $16$, l'enter $298$ ($12A$ en hexadecimal) només és $2$-espaiat.
\end{itemize}

\medskip

Donada una seqüència de casos, a on cada cas està format per un enter $b \geq 2$, un enter $0 \leq d \leq b$, un enter $n \geq 0$, i una seqüència d'$n$ enters no negatius volem saber, per cada cas, quants elements de la seqüència són $d$-espaiats en base $b$.

\Input

L'entrada és una seqüència de casos. Cada cas comença amb una base $b \geq 2$, una xifra $0 \leq d \leq b$ i un enter no negatiu $n$, i li segueix una seqüència d'$n$ enters no negatius.


\Output

Per cada cas, s'ha de dir quants dels elements de la seqüència d'enters és $d$-espaiat en base $b$.

\Sample
